\section{A tétel lényege és vizsgaszintű vázlata}

\begin{summarybox}
A 20. tétel az SQL nyelvet mint relációs adatbázis-kezelő nyelvet tárgyalja. Az SQL egyszerre tartalmaz szerkezetkezelő, adatkezelő, lekérdező és jogosultságkezelő utasításokat. A DDL az adatbázis-objektumok létrehozásával, módosításával és törlésével foglalkozik; a DML a rekordok beszúrását, módosítását és törlését adja; a DCL a jogosultságok kezelésére szolgál; a \code{SELECT} pedig az adatok lekérdezésének központi utasítása. A tételben különösen fontos a relációs algebra és az SQL kapcsolata: a szelekció, projekció, join, aggregáció és halmazműveletek mind megfogalmazhatók SQL-ben.
\end{summarybox}

Vizsgán célszerű a következő sorrendet követni:

\begin{enumerate}
  \item SQL szerepe: magas szintű, deklaratív, relációs adatbázis-kezelő nyelv;
  \item SQL parancscsoportok: DDL, DML, DQL/\code{SELECT}, DCL, kiegészítésként TCL;
  \item DDL: \code{CREATE TABLE}, adattípusok, integritási feltételek, \code{ALTER TABLE}, \code{DROP TABLE};
  \item DML: \code{INSERT}, \code{UPDATE}, \code{DELETE};
  \item DCL: jogosultság, privilégium, szerepkör, \code{GRANT}, \code{REVOKE};
  \item \code{SELECT}: \code{FROM}, \code{JOIN}, \code{WHERE}, \code{GROUP BY}, \code{HAVING}, \code{SELECT}, \code{ORDER BY};
  \item beágyazott \code{SELECT}: alselect a \code{WHERE}, \code{FROM} és \code{SELECT} részekben, korrelált és nem korrelált alselect;
  \item relációs algebra SQL-ben: szelekció, projekció, kiterjesztés, joinok, aggregáció, halmazműveletek, osztás;
  \item gyakori vizsgahibák és rövid, összefüggő válaszminta.
\end{enumerate}

A legfontosabb vizsgamondat: az SQL deklaratív adatbázis-nyelv, amelyben a DDL a sémát, a DML az adatokat, a DCL a jogosultságokat kezeli, a \code{SELECT} pedig relációs algebrai műveletekre építve állít elő eredményrelációt a megadott táblákból.

\begin{center}
\begin{tabularx}{\textwidth}{L{22mm}L{31mm}X}
\toprule
\textbf{Csoport} & \textbf{Tipikus utasítások} & \textbf{Feladat} \\
\midrule
DDL & \code{CREATE}, \code{ALTER}, \code{DROP} & Adatbázis-objektumok, például táblák, nézetek, indexek és sémák létrehozása, módosítása és megszüntetése. \\
DML & \code{INSERT}, \code{UPDATE}, \code{DELETE} & A táblákban tárolt rekordok felvitele, módosítása és törlése. \\
DQL & \code{SELECT} & Adatok lekérdezése, szűrése, összekapcsolása, csoportosítása és számított eredmények előállítása. \\
DCL & \code{GRANT}, \code{REVOKE} & Jogosultságok, privilégiumok és szerepkörök kiosztása vagy visszavonása. \\
TCL & \code{COMMIT}, \code{ROLLBACK}, \code{SAVEPOINT} & Tranzakciók véglegesítése, visszavonása és részleges kezelése; sok tananyagban külön kiegészítő csoportként szerepel. \\
\bottomrule
\end{tabularx}

\end{center}

\section{Az SQL szerepe és szemlélete}

\subsection{Mi az SQL?}

Az \textbf{SQL} a Structured Query Language rövidítése. Relációs adatbázis-kezelő rendszerekben használatos magas szintű nyelv, amelynek segítségével adatbázis-objektumokat hozhatunk létre, adatokat módosíthatunk, adatokat kérdezhetünk le és jogosultságokat kezelhetünk.

Az SQL szemlélete alapvetően \textbf{deklaratív}. A felhasználó többnyire azt írja le, hogy milyen eredményt szeretne látni, nem azt, hogy az adatbázis-kezelő pontosan milyen algoritmussal járja be a táblákat. Például egy \code{SELECT} utasításban megadjuk a forrástáblákat, a kapcsolási és szűrési feltételeket, valamint az eredmény mezőit. A fizikai végrehajtási tervet az adatbázis-kezelő optimalizálója választja ki.

Ez fontos különbség az algoritmikus programozási nyelvekhez képest. Egy C, Java vagy C\# programban gyakran ciklusokkal és lépésenkénti feldolgozással írjuk le a teendőt. SQL-ben inkább halmazokban gondolkodunk: táblákból új eredménytáblát állítunk elő.

\subsection{SQL szabvány és dialektusok}

Az SQL szabványosított nyelv, de a gyakorlatban az egyes adatbázis-kezelők saját kiterjesztéseket és eltérő részleteket használnak. Oracle, PostgreSQL, Microsoft SQL Server, MySQL/MariaDB és SQLite alatt az alapvető \code{SELECT}, \code{INSERT}, \code{UPDATE}, \code{DELETE}, \code{CREATE TABLE} gondolkodás hasonló, de eltérhetnek például az adattípusok nevei, az automatikus azonosítók kezelése, a dátumfüggvények, a jogosultságkezelés és egyes halmazműveletek nevei.

Vizsgán ezért érdemes az SQL-t szabványos, általános szinten bemutatni, és jelezni, ha egy részlet DBMS-függő. Például a különbségműveletet több rendszer \code{EXCEPT} néven, Oracle környezetben gyakran \code{MINUS} néven támogatja.

\subsection{SQL és a relációs modell kapcsolata}

A relációs modellben az adatok relációkban, vagyis táblákhoz hasonló szerkezetekben jelennek meg. Az SQL erre a modellre épít, de gyakorlati adatbázis-nyelvként több ponton eltér az elméleti relációs algebrától.

Fontos különbségek:

\begin{itemize}
  \item a relációs algebra elméletileg halmazokkal dolgozik, SQL-ben viszont ismétlődő sorok is megjelenhetnek, hacsak nem használunk \code{DISTINCT} kulcsszót;
  \item a relációs modellben a soroknak nincs sorrendjük, SQL-ben az eredmény megjelenítési sorrendjét \code{ORDER BY} záradékkal kérhetjük;
  \item SQL-ben létezik \code{NULL}, vagyis hiányzó/ismeretlen érték, amely háromértékű logikai viselkedést okozhat;
  \item SQL-ben számos gyakorlati kiterjesztés van, például nézetek, indexek, tárolt eljárások és DBMS-specifikus függvények.
\end{itemize}

A relációs algebra mégis nagyon fontos, mert megadja a lekérdezések logikai alapját. A \code{WHERE} a szelekcióhoz, a \code{SELECT} mezőlista a projekcióhoz, a \code{JOIN} a relációk összekapcsolásához, a \code{GROUP BY} pedig a csoportképzéshez kapcsolható.

\section{SQL parancscsoportok}

\subsection{DDL: adatdefiníciós nyelv}

A \textbf{DDL}, vagyis Data Definition Language az adatbázis szerkezetét kezeli. Ide tartoznak azok az utasítások, amelyek adatbázis-objektumokat hoznak létre, módosítanak vagy törölnek.

Tipikus DDL utasítások:

\begin{itemize}
  \item \code{CREATE}: objektum létrehozása, például \code{CREATE TABLE}, \code{CREATE VIEW}, \code{CREATE INDEX};
  \item \code{ALTER}: meglévő objektum szerkezetének módosítása, például új oszlop vagy új megkötés felvétele;
  \item \code{DROP}: objektum megszüntetése.
\end{itemize}

DDL-nél a fő gondolat az, hogy nem konkrét rekordokat kezelünk, hanem a sémát. Egy tábla létrehozásakor megadjuk a tábla nevét, mezőit, azok típusát és a rájuk vonatkozó integritási feltételeket.

\subsection{DML: adatkezelő nyelv}

A \textbf{DML}, vagyis Data Manipulation Language a táblák tartalmát kezeli. A DML utasítások már konkrét rekordokra hatnak.

Tipikus DML utasítások:

\begin{itemize}
  \item \code{INSERT}: új rekord vagy rekordok beszúrása;
  \item \code{UPDATE}: meglévő rekordok mezőértékeinek módosítása;
  \item \code{DELETE}: rekordok törlése.
\end{itemize}

A DML utasításoknál különösen fontos a \code{WHERE} feltétel. Ha egy \code{UPDATE} vagy \code{DELETE} utasításból kimarad a szűrés, akkor a művelet akár a tábla összes rekordjára is vonatkozhat.

\subsection{DQL: lekérdezés}

Sok felosztásban a \code{SELECT} önálló csoportként, \textbf{DQL} néven szerepel. A \code{SELECT} nem módosítja közvetlenül a tárolt adatokat, hanem eredményrelációt állít elő. Ezzel lehet megvalósítani a relációs algebra lekérdező műveleteit.

A \code{SELECT} egyszerű alakja:

\begin{center}
\code{SELECT mezők FROM tábla WHERE feltétel;}
\end{center}

Összetettebb lekérdezésben joinok, csoportosítás, aggregáció, alselectek, halmazműveletek és rendezés is szerepelhetnek.

\subsection{DCL: jogosultságkezelés}

A \textbf{DCL}, vagyis Data Control Language az adatbázis-védelmi és jogosultsági műveleteket fedi le. A fő fogalmak:

\begin{itemize}
  \item \textbf{felhasználó}: adatbázisba belépő szubjektum;
  \item \textbf{objektum}: például tábla, nézet, eljárás vagy séma;
  \item \textbf{privilégium}: valamilyen művelet joga, például \code{SELECT}, \code{INSERT}, \code{UPDATE};
  \item \textbf{szerepkör}: jogosultságok halmaza, amelyet több felhasználóhoz is hozzá lehet rendelni.
\end{itemize}

A DCL legfontosabb utasításai:

\begin{itemize}
  \item \code{GRANT}: jogosultság adása;
  \item \code{REVOKE}: jogosultság visszavonása.
\end{itemize}

\section{DDL: táblák és adatbázis-objektumok kezelése}

\subsection{Adatbázis-objektumok}

SQL-ben objektumnak tekinthető például a tábla, nézet, index, séma, tárolt eljárás, függvény, trigger, sequence, felhasználó vagy szerepkör. A 20. tétel szempontjából a legfontosabb objektum a \textbf{tábla}, mert a relációs adatmodell gyakorlati megvalósításában az adatokat táblákban tároljuk.

Egy tábla definíciója tartalmazza:

\begin{itemize}
  \item a tábla nevét;
  \item az oszlopok nevét;
  \item az oszlopok adattípusát;
  \item mezőszintű és táblaszintű integritási feltételeket;
  \item esetenként DBMS-függő tárolási vagy indexelési paramétereket.
\end{itemize}

\subsection{CREATE TABLE}

A \code{CREATE TABLE} utasítás új táblát hoz létre. Általános alakja:

\begin{center}
\code{CREATE TABLE táblanév (oszlopnév típus megkötések, ...);}
\end{center}

A mezők típusai DBMS-függők lehetnek, de tipikus kategóriák:

\begin{itemize}
  \item egész és valós számok: például \code{INTEGER}, \code{NUMERIC}, \code{DECIMAL}, \code{FLOAT};
  \item karakteres adatok: például \code{CHAR}, \code{VARCHAR};
  \item dátum és idő: például \code{DATE}, \code{TIME}, \code{TIMESTAMP};
  \item nagy objektumok: például \code{BLOB}, \code{CLOB}.
\end{itemize}

A táblalétrehozásnál már érdemes megadni a kulcsokat és a legfontosabb integritási feltételeket, mert ezek védik az adatbázis helyességét.

\begin{center}
\begin{tabularx}{\textwidth}{L{33mm}L{37mm}X}
\toprule
\textbf{Megkötés} & \textbf{Példa} & \textbf{Jelentés} \\
\midrule
Elsődleges kulcs & \code{PRIMARY KEY} & Egyedileg azonosítja a rekordokat; nem lehet \code{NULL}. \\
Idegen kulcs & \code{FOREIGN KEY ... REFERENCES} & Másik tábla kulcsára hivatkozik, ezzel hivatkozási integritást biztosít. \\
Nem üres érték & \code{NOT NULL} & A mezőnek minden rekordban értéket kell kapnia. \\
Egyediség & \code{UNIQUE} & A mező vagy mezőcsoport értékei nem ismétlődhetnek. \\
Értékellenőrzés & \code{CHECK (ar > 0)} & Logikai feltétellel korlátozza a megengedett mező- vagy rekordértékeket. \\
Alapértelmezés & \code{DEFAULT 0} & Ha beszúráskor nem adunk értéket, a rendszer az alapértelmezett értéket használja. \\
\bottomrule
\end{tabularx}

\end{center}

\zvcode[caption={Egyszerű táblaséma DDL utasításokkal}]{src/code/zv20_schema.sql}

A példában a \code{Gyarto} tábla azonosítója elsődleges kulcs. A \code{Termek} tábla \code{gyartoId} mezője idegen kulcs, amely a \code{Gyarto} táblára hivatkozik. Ez azt jelenti, hogy nem lehet olyan terméket beszúrni, amely nem létező gyártóra mutat.

\subsection{Mezőszintű és táblaszintű megkötések}

A megkötések két gyakori formában jelenhetnek meg:

\begin{itemize}
  \item \textbf{mezőszintű megkötés}: közvetlenül az oszlop definíciója mellett szerepel, például \code{ar INTEGER CHECK (ar > 0)};
  \item \textbf{táblaszintű megkötés}: külön sorban, több oszlopra is vonatkozhat, például összetett elsődleges kulcs vagy összetett idegen kulcs esetén.
\end{itemize}

Összetett kulcsnál a táblaszintű forma természetesebb:

\begin{center}
\code{PRIMARY KEY (hallgatoId, targyId)}
\end{center}

Idegen kulcsnál fontos, hogy a hivatkozott tábla és a hivatkozott kulcs létezzen. Törlésnél és módosításnál a hivatkozási integritás miatt az adatbázis-kezelő tiltást, kaszkád műveletet, \code{NULL}-ra állítást vagy alapértelmezett értékre állítást alkalmazhat, attól függően, hogyan definiáltuk a kapcsolatot.

\subsection{ALTER TABLE és DROP TABLE}

Az \code{ALTER TABLE} meglévő tábla szerkezetét módosítja. Tipikus használata:

\begin{itemize}
  \item új oszlop hozzáadása;
  \item oszlop típusának vagy megkötésének módosítása;
  \item új kulcs vagy \code{CHECK} feltétel felvétele;
  \item oszlop vagy megkötés törlése, ha ezt az adatbázis-kezelő és az aktuális adatok megengedik.
\end{itemize}

A \code{DROP TABLE} megszünteti a táblát. Ezzel a tábla szerkezete és tartalma is eltűnik. Ha más tábla idegen kulccsal hivatkozik rá, a törlés gyakran nem hajtható végre addig, amíg a hivatkozásokat nem kezeljük.

\begin{warningbox}
A DDL utasításoknál mindig figyelni kell a függőségekre. Egy hivatkozott táblát, kulcsot vagy oszlopot nem lehet következmény nélkül törölni, ha más objektumok épülnek rá.
\end{warningbox}

\section{DML: rekordok felvitele, módosítása és törlése}

\subsection{INSERT}

Az \code{INSERT} új rekordot visz fel egy táblába. Két gyakori alakja:

\begin{itemize}
  \item \code{INSERT INTO tábla VALUES (...)}: ilyenkor az értéklista sorrendjének és darabszámának illeszkednie kell a tábla oszlopsorrendjéhez;
  \item \code{INSERT INTO tábla (mezőlista) VALUES (...)}: biztonságosabb és áttekinthetőbb, mert megmondja, mely oszlopok kapnak értéket.
\end{itemize}

Ha egy mező nem kap értéket, akkor az adatbázis-kezelő vagy az alapértelmezett értéket használja, vagy \code{NULL}-t állít be, ha ez megengedett. Ha a mező \code{NOT NULL} és nincs alapértelmezett értéke, a beszúrás hibát ad.

\subsection{UPDATE}

Az \code{UPDATE} meglévő rekordok mezőértékeit módosítja. A \code{SET} rész adja meg az új értékeket, a \code{WHERE} rész pedig azt, hogy mely rekordokra vonatkozik a módosítás.

Általános alak:

\begin{center}
\code{UPDATE tábla SET mező = érték WHERE feltétel;}
\end{center}

Az \code{UPDATE} is integritási ellenőrzések alá esik. Nem állíthatunk be olyan idegen kulcsértéket, amely nem létező rekordra hivatkozik, és nem sérthetjük meg a \code{CHECK}, \code{UNIQUE} vagy \code{NOT NULL} feltételeket sem.

\subsection{DELETE}

A \code{DELETE} rekordokat töröl egy táblából. Általános alak:

\begin{center}
\code{DELETE FROM tábla WHERE feltétel;}
\end{center}

Ha a \code{WHERE} hiányzik, a művelet a tábla minden rekordját törölheti. Ez vizsgán is gyakori hibapont: a \code{DELETE FROM Termek;} nem egy terméket töröl, hanem minden terméket, ha az adatbázis-kezelő és az integritási feltételek ezt engedik.

\zvcode[caption={DML és DCL alaputasítások példái}]{src/code/zv20_dml_dcl.sql}

\section{DCL: jogosultságok és szerepkörök}

\subsection{Jogosultságkezelés célja}

Egy adatbázist több felhasználó vagy alkalmazás is használhat. Nem mindenkinek szabad ugyanazt látnia vagy módosítania. A jogosultságkezelés feladata, hogy szabályozza, ki milyen objektumon milyen műveletet végezhet.

Például:

\begin{itemize}
  \item egy ügyintéző olvashat és beszúrhat rendeléseket;
  \item egy adminisztrátor módosíthat sémát és jogosultságot adhat;
  \item egy publikus felhasználó csak nézeten keresztül láthat szűrt adatokat;
  \item egy alkalmazás csak a működéséhez szükséges minimális jogokat kapja meg.
\end{itemize}

A minimális jogosultság elve szerint minden felhasználónak és alkalmazásnak csak azokat a jogokat érdemes megadni, amelyek tényleg szükségesek a feladatához.

\subsection{GRANT és REVOKE}

A \code{GRANT} jogosultságot ad:

\begin{center}
\code{GRANT SELECT ON Termek TO ugyintezo;}
\end{center}

A \code{REVOKE} jogosultságot von vissza:

\begin{center}
\code{REVOKE SELECT ON Termek FROM ugyintezo;}
\end{center}

Szerepkör használatával nem egyesével kell minden felhasználónak minden jogot kiosztani. Létrehozható például egy \code{lekerdezo} szerepkör, amely megkapja a szükséges \code{SELECT} jogokat, majd ezt a szerepkört lehet a felhasználókhoz rendelni. A szerepkörök karbantarthatóbbá teszik a jogosultságkezelést.

\section{A SELECT utasítás}

\subsection{A SELECT alapgondolata}

A \code{SELECT} az SQL lekérdező utasítása. Egy vagy több forrásból eredményrelációt állít elő. Egyszerű esetben egy táblából kiválaszt néhány oszlopot és néhány sort:

\begin{center}
\code{SELECT nev, ar FROM Termek WHERE ar > 1000;}
\end{center}

Ez relációs algebrai szemmel projekció és szelekció együtt: a \code{WHERE} kiválasztja a megfelelő rekordokat, a \code{SELECT nev, ar} pedig csak a szükséges attribútumokat tartja meg.

\begin{figure}[htbp]
\centering
\resizebox{0.96\textwidth}{!}{%
\begin{tikzpicture}[font=\scriptsize, node distance=9mm]
  \node[zvflowprocess, fill=black!5, minimum width=36mm] (from) {\code{FROM}\ és \code{JOIN}\\források és kapcsolás};
  \node[zvflowprocess, fill=lightaccent, right=of from, minimum width=31mm] (where) {\code{WHERE}\\rekordok szűrése};
  \node[zvflowprocess, fill=black!5, right=of where, minimum width=33mm] (group) {\code{GROUP BY}\\csoportképzés};
  \node[zvflowprocess, fill=lightaccent, right=of group, minimum width=31mm] (having) {\code{HAVING}\\csoportok szűrése};
  \node[zvflowprocess, fill=warnbg, below=of group, minimum width=35mm] (select) {\code{SELECT}\\projekció, számítás};
  \node[zvflowprocess, fill=black!5, left=of select, minimum width=31mm] (distinct) {\code{DISTINCT}\\ismétlődések elhagyása};
  \node[zvflowprocess, fill=lightaccent, left=of distinct, minimum width=31mm] (order) {\code{ORDER BY}\\megjelenítési sorrend};

  \draw[arrow] (from) -- (where);
  \draw[arrow] (where) -- (group);
  \draw[arrow] (group) -- (having);
  \draw[arrow] (having) |- (select);
  \draw[arrow] (select) -- (distinct);
  \draw[arrow] (distinct) -- (order);
\end{tikzpicture}%
}
\caption{A \code{SELECT} lekérdezés tipikus logikai feldolgozási sorrendje}
\label{fig:zv20_select_logikai_sorrend}
\end{figure}


A logikai sorrend nem mindig azonos a fizikai végrehajtási sorrenddel. Az optimalizáló átrendezheti a műveleteket, például korábban elvégezhet szűréseket vagy indexet használhat. A logikai sorrend mégis segít megérteni, hogy miért nem használhatunk minden záradékban minden aliasnevet vagy aggregált kifejezést.

\begin{center}
\begin{tabularx}{\textwidth}{L{25mm}L{32mm}X}
\toprule
\textbf{Záradék} & \textbf{Relációs algebrai szerep} & \textbf{Feladat} \\
\midrule
\code{FROM} & Operandusrelációk kijelölése & Megadja, mely táblákból, nézetekből vagy részlekérdezésekből indul a lekérdezés. \\
\code{JOIN ... ON} & Join, Descartes-szorzat szűréssel & Összekapcsolja a forrásokat, gyakran elsődleges kulcs--idegen kulcs kapcsolat alapján. \\
\code{WHERE} & Szelekció & Rekordszintű feltétellel szűri az eredményt a csoportosítás előtt. \\
\code{GROUP BY} & Csoportképzés & A rekordokat csoportokra bontja, hogy csoportonként lehessen aggregálni. \\
\code{HAVING} & Csoportokra vonatkozó szelekció & Aggregált értékek alapján szűri a csoportokat. \\
\code{SELECT} & Projekció és kiterjesztés & Meghatározza, mely mezők és számított kifejezések kerüljenek az eredménybe. \\
\code{ORDER BY} & Rendezés, nem tisztán algebrai művelet & Megadja az eredmény megjelenítési sorrendjét; a relációs algebra halmazszemlélete miatt ez külön SQL-funkció. \\
\bottomrule
\end{tabularx}

\end{center}

\subsection{Projekció és kiterjesztés}

A \code{SELECT} mezőlistája adja meg, hogy milyen oszlopok kerüljenek az eredménybe. Ha minden oszlop kell, használható a \code{*}, de vizsgán és valós fejlesztésben jobb név szerint felsorolni a szükséges mezőket.

Példák:

\begin{itemize}
  \item \code{SELECT nev, ar FROM Termek}: csak a név és ár oszlop jelenik meg;
  \item \code{SELECT nev, ar * 1.27 AS bruttoAr FROM Termek}: számított oszlopot hoz létre;
  \item \code{SELECT DISTINCT kategoria FROM Termek}: ismétlődésmentesen listázza a kategóriákat.
\end{itemize}

A \code{DISTINCT} azért fontos, mert SQL-ben az eredmény alapértelmezés szerint nem feltétlenül halmaz, hanem többszörösen előforduló sorokat is tartalmazhat.

\subsection{Szelekció WHERE feltétellel}

A \code{WHERE} rekordszintű feltételt ad. Csak azok a rekordok jutnak tovább, amelyekre a feltétel igaz.

Gyakori feltételtípusok:

\begin{itemize}
  \item összehasonlítás: \code{ar > 1000}, \code{nev = 'Tej'};
  \item logikai kapcsolás: \code{AND}, \code{OR}, \code{NOT};
  \item mintaillesztés: \code{LIKE 'A\%'};
  \item intervallum: \code{BETWEEN 100 AND 500};
  \item halmaztagság: \code{IN (...)};
  \item \code{NULL} vizsgálat: \code{IS NULL}, \code{IS NOT NULL}.
\end{itemize}

A \code{NULL} értékkel külön kell bánni. A \code{telephely = NULL} nem helyes SQL-feltétel; helyette \code{telephely IS NULL} szükséges.

\subsection{JOIN műveletek}

A \code{JOIN} több táblát kapcsol össze. A relációs adatmodellben a kapcsolatok gyakran idegen kulccsal jelennek meg, a lekérdezésben pedig joinnal kapcsoljuk vissza az összetartozó rekordokat.

Gyakori joinok:

\begin{itemize}
  \item \textbf{inner join}: csak az illeszkedő rekordpárok jelennek meg;
  \item \textbf{left outer join}: a bal oldali tábla minden rekordja megmarad, a pár nélküli jobb oldali értékek \code{NULL}-ok lesznek;
  \item \textbf{right outer join}: a jobb oldali tábla minden rekordja megmarad;
  \item \textbf{full outer join}: mindkét oldal pár nélküli rekordjai is megmaradnak;
  \item \textbf{cross join}: Descartes-szorzatot képez.
\end{itemize}

A joinfeltételt célszerű \code{ON} záradékban megadni:

\begin{center}
\code{FROM Termek t INNER JOIN Gyarto g ON t.gyartoId = g.gyartoId}
\end{center}

Az aliasok, például \code{t} és \code{g}, rövidebbé és egyértelműbbé teszik a lekérdezést.

\subsection{Aggregáció, GROUP BY és HAVING}

Az aggregáció több rekordból számít egy összesített értéket. Tipikus aggregáló függvények:

\begin{itemize}
  \item \code{COUNT(*)}: rekordok száma;
  \item \code{SUM(mező)}: összeg;
  \item \code{AVG(mező)}: átlag;
  \item \code{MIN(mező)}, \code{MAX(mező)}: minimum és maximum.
\end{itemize}

Ha az egész táblára kérünk összesítést, nem kell \code{GROUP BY}. Ha csoportonként szeretnénk összesítést, meg kell adni a csoportképző mezőket.

Példa: gyártónként termékszám és átlagár:

\begin{center}
\code{GROUP BY g.gyartoId, g.nev}
\end{center}

A \code{WHERE} a rekordokat szűri csoportosítás előtt. A \code{HAVING} a már létrejött csoportokat szűri, ezért aggregált kifejezésekre is hivatkozhat, például \code{HAVING COUNT(*) >= 2}.

\zvcode[caption={SELECT példák szűréssel, joinnal, aggregációval és alselecttel}]{src/code/zv20_select_peldak.sql}

\subsection{ORDER BY}

A relációs algebra halmazszemléletű, ezért az eredmény sorrendje nem része a reláció jelentésének. SQL-ben a megjelenített eredmény sorrendjét \code{ORDER BY} adja meg.

Példák:

\begin{itemize}
  \item \code{ORDER BY nev}: név szerint növekvő sorrend;
  \item \code{ORDER BY ar DESC}: ár szerint csökkenő sorrend;
  \item \code{ORDER BY kategoria, ar DESC}: kategórián belül ár szerint csökkenő sorrend.
\end{itemize}

\section{Beágyazott SELECT használata}

\subsection{Az alselect fogalma}

A beágyazott \code{SELECT}, más néven alselect vagy részlekérdezés olyan lekérdezés, amely egy másik SQL utasításon belül szerepel. Részlekérdezést használhatunk feltételben, forrásként, számított mezőként vagy DML utasítás részeként is.

\begin{center}
\begin{tabularx}{\textwidth}{L{31mm}L{35mm}X}
\toprule
\textbf{Típus} & \textbf{Elhelyezés} & \textbf{Jellemző használat} \\
\midrule
Skalár alselect & \code{SELECT} listában vagy feltételben & Egyetlen értéket ad vissza, például átlagárat vagy darabszámot. \\
Táblaszerű alselect & \code{FROM} záradékban & Ideiglenes, névvel ellátott forrásként használható, amelyre külső lekérdezés épül. \\
Nem korrelált alselect & Külső sorra nem hivatkozik & Egyszer önállóan kiértékelhető; például \code{IN (SELECT ...)}. \\
Korrelált alselect & Belső lekérdezés hivatkozik a külső sorra & A belső lekérdezés a külső rekord aktuális értékétől függ; gyakori \code{EXISTS} mellett. \\
Halmaztesztes alselect & \code{IN}, \code{EXISTS}, \code{ANY}, \code{ALL} & A külső rekordot részlekérdezés eredményhalmazához hasonlítja. \\
\bottomrule
\end{tabularx}

\end{center}

Az alselect erős eszköz, mert egy lekérdezés eredményét közvetlenül felhasználhatjuk egy másik lekérdezésben. Ugyanakkor figyelni kell arra, hogy a részlekérdezés hány sort és hány oszlopot ad vissza. Egy skalár összehasonlítás egyetlen értéket vár, az \code{IN} viszont több értéket tartalmazó halmazt is elfogad.

\subsection{Alselect a WHERE feltételben}

A leggyakoribb forma, amikor a részlekérdezés a \code{WHERE} feltételben szerepel.

Példa: az átlagárnál drágább termékek:

\begin{center}
\code{WHERE ar > (SELECT AVG(ar) FROM Termek)}
\end{center}

Példa: azok a termékek, amelyek egy adott telephelyű gyártóhoz tartoznak:

\begin{center}
\code{WHERE gyartoId IN (SELECT gyartoId FROM Gyarto WHERE telephely = 'Miskolc')}
\end{center}

A \code{NOT IN} használatánál óvatosnak kell lenni, ha a részlekérdezés \code{NULL} értéket adhat vissza. Ilyenkor sokszor biztonságosabb a \code{NOT EXISTS} megoldás.

\subsection{Korrelált alselect és EXISTS}

Korrelált alselect esetén a belső lekérdezés hivatkozik a külső lekérdezés aktuális sorára. Emiatt a belső lekérdezés logikailag minden külső rekordhoz kapcsolódva értelmezhető.

Példa: azok a gyártók, amelyeknek nincs termékük:

\begin{center}
\code{WHERE NOT EXISTS (SELECT 1 FROM Termek t WHERE t.gyartoId = g.gyartoId)}
\end{center}

Az \code{EXISTS} nem a belső lekérdezés konkrét oszlopértékét használja, hanem azt vizsgálja, hogy létezik-e legalább egy illeszkedő rekord.

\subsection{Alselect a FROM és SELECT részekben}

A \code{FROM} részben szereplő alselect táblaszerű eredményt ad, amelyet aliasnévvel láthatunk el. Ez hasznos, ha először összesíteni szeretnénk, majd az összesített eredményre további feltételt vagy join műveletet alkalmazunk.

A \code{SELECT} listában szereplő alselect tipikusan skalár értéket ad. Például minden termék mellé kiírhatjuk az összes termék átlagárát. Ilyenkor figyelni kell arra, hogy a belső lekérdezés tényleg egyetlen értéket adjon vissza.

\section{Relációs algebra SQL-ben}

\subsection{Áttekintés}

A relációs algebra absztrakt, logikai műveleteket ad relációkon. SQL-ben ezek gyakorlati parancsformában jelennek meg. A megfeleltetés nem mindig teljesen egy az egyben történik, mert SQL-ben van \code{NULL}, ismétlődés, rendezés, alias, függvények és DBMS-függő kiterjesztések is.

\begin{center}
\begin{tabularx}{\textwidth}{L{25mm}L{43mm}X}
\toprule
\textbf{Relációs algebra} & \textbf{SQL-megfelelő} & \textbf{Megjegyzés} \\
\midrule
Szelekció $\sigma_{felt}(R)$ & \code{WHERE felt} & Rekordokat szűr; a séma alapvetően nem változik. \\
Projekció $\pi_{A,B}(R)$ & \code{SELECT A, B} & Attribútumokat választ ki; SQL-ben \code{DISTINCT} kell, ha halmazszerű ismétlődésmentes eredményt akarunk. \\
Kiterjesztés & \code{SELECT kifejezés AS név} & Számított mezőt hoz létre, például árkedvezményt vagy kategóriát. \\
Descartes-szorzat $R \times S$ & \code{CROSS JOIN} & Minden rekordpárt előállít; ritkán kell közvetlenül, mert gyorsan nagy eredményt ad. \\
Szelekciós join $R \bowtie_{felt} S$ & \code{INNER JOIN ... ON felt} & A kapcsolódó rekordpárokat tartja meg. \\
Külső join & \code{OUTER JOIN} & A nem illeszkedő rekordokat is megtartja az egyik vagy mindkét oldalról. \\
Aggregáció és csoportképzés & \code{GROUP BY}, \code{COUNT}, \code{SUM}, \code{AVG} & Csoportonként összesítő értékeket számít. \\
Unió & \code{UNION} vagy \code{UNION ALL} & \code{UNION} ismétlődésmentesít, \code{UNION ALL} megtartja az ismétlődéseket. \\
Metszet & \code{INTERSECT} vagy joinos megoldás & Kompatibilis sémájú lekérdezések közös rekordjait adja. \\
Különbség & \code{EXCEPT} / \code{MINUS} vagy \code{NOT EXISTS} & DBMS-függő névhasználat lehet; Oracle-ben gyakori a \code{MINUS}. \\
Osztás & \code{NOT EXISTS}-es kettős tagadás & Olyan elemeket keres, amelyek minden keresett értékkel kapcsolatban állnak. \\
\bottomrule
\end{tabularx}

\end{center}

\subsection{Szelekció és projekció SQL-ben}

A szelekció rekordokat választ ki. SQL-ben ezt a \code{WHERE} végzi:

\begin{center}
\code{SELECT * FROM Termek WHERE ar > 1000;}
\end{center}

A projekció attribútumokat választ ki. SQL-ben ez a \code{SELECT} mezőlistája:

\begin{center}
\code{SELECT nev, ar FROM Termek;}
\end{center}

A két művelet természetesen kombinálható:

\begin{center}
\code{SELECT nev, ar FROM Termek WHERE ar > 1000;}
\end{center}

Relációs algebrai alakban ez:

\[
  \pi_{nev,ar}(\sigma_{ar>1000}(Termek)).
\]

\subsection{Join SQL-ben}

A join több reláció rekordjait kapcsolja össze. SQL-ben ezt explicit \code{JOIN} szintaxissal érdemes írni:

\begin{center}
\code{SELECT t.nev, g.nev FROM Termek t INNER JOIN Gyarto g ON t.gyartoId = g.gyartoId;}
\end{center}

Relációs algebrai szemmel ez egy szelekciós join. A feltétel általában kulcs--idegen kulcs egyezés. Ha a gyártó nélküli vagy termék nélküli rekordokat is látni szeretnénk, külső joint kell választani.

\subsection{Aggregáció és csoportképzés SQL-ben}

Az aggregáció relációs algebrai jelölése gyakran \(\gamma\). SQL-ben az aggregáló függvények és a \code{GROUP BY} valósítják meg.

Példa:

\begin{center}
\code{SELECT gyartoId, COUNT(*) FROM Termek GROUP BY gyartoId;}
\end{center}

Ha a csoportokra feltételt is adunk, \code{HAVING} szükséges:

\begin{center}
\code{HAVING COUNT(*) > 3}
\end{center}

Fontos szabály: ha \code{GROUP BY} szerepel, akkor a \code{SELECT} listában lévő nem aggregált mezőknek általában szerepelniük kell a \code{GROUP BY} részben is.

\subsection{Halmazműveletek SQL-ben}

Az SQL támogatja a lekérdezések eredményhalmazainak kombinálását. A forráslekérdezéseknek kompatibilis oszlopszámmal és összeegyeztethető típusokkal kell rendelkezniük.

\begin{itemize}
  \item \code{UNION}: két eredmény uniója ismétlődésmentesítéssel;
  \item \code{UNION ALL}: unió ismétlődések megtartásával;
  \item \code{INTERSECT}: közös sorok;
  \item \code{EXCEPT} vagy \code{MINUS}: az első eredményből elhagyja a másodikban is szereplő sorokat.
\end{itemize}

A \code{UNION ALL} gyakran gyorsabb lehet, mert nem kell ismétlődéseket eltávolítani. Ha viszont relációs algebrai értelemben halmazt szeretnénk, akkor ismétlődésmentes eredmény kell.

\subsection{Osztás művelet SQL-ben}

A relációs algebra osztás művelete olyan kérdésekre használható, ahol „minden” típusú feltétel szerepel. Például: mely gyártók gyártanak minden keresett kategóriából terméket?

SQL-ben nincs minden rendszerben közvetlen \code{DIVIDE} utasítás, ezért gyakran \code{NOT EXISTS}-es kettős tagadással írjuk fel. A gondolat:

\begin{itemize}
  \item válasszuk ki azt a gyártót;
  \item amelyhez nem létezik olyan keresett kategória;
  \item amelyhez nem létezik megfelelő termék.
\end{itemize}

Ez elsőre furcsa, de logikailag azt jelenti: minden keresett kategóriához van illeszkedő terméke.

\zvcode[caption={Relációs algebrai műveletek SQL-megfelelői példákkal}]{src/code/zv20_relacios_algebra_sql.sql}

\section{Összefüggő példa}

\subsection{Példaséma}

Legyen két táblánk:

\[
  Gyarto(gyartoId, nev, telephely)
\]

\[
  Termek(termekId, nev, ar, kategoria, gyartoId, keszlet)
\]

A \code{Termek.gyartoId} idegen kulcs a \code{Gyarto.gyartoId} mezőre. Ez a séma alkalmas arra, hogy bemutassuk a DDL, DML és \code{SELECT} alapjait.

\subsection{Műveletek a példasémán}

DDL szinten a \code{CREATE TABLE} létrehozza a két táblát, elsődleges kulccsal, idegen kulccsal és értékellenőrzéssel. DML szinten az \code{INSERT} gyártót és terméket szúr be, az \code{UPDATE} módosítja az árakat, a \code{DELETE} törli a felesleges rekordokat. DCL szinten a \code{GRANT} olvasási vagy módosítási jogot adhat a megfelelő felhasználónak.

Lekérdezési szinten:

\begin{itemize}
  \item a drágább termékek kiválasztása \code{WHERE} feltétellel szelekció;
  \item a terméknév és ár kiírása projekció;
  \item a termékek és gyártók összekapcsolása join;
  \item a gyártónkénti termékszám \code{GROUP BY} és \code{COUNT};
  \item az átlagárnál drágább termékek keresése beágyazott \code{SELECT};
  \item a termékkel nem rendelkező gyártók keresése \code{NOT EXISTS} vagy külső join segítségével oldható meg.
\end{itemize}

\subsection{Relációs algebrai és SQL gondolkodás összevetése}

Relációs algebrai alak:

\[
  \pi_{g.nev,t.nev,t.ar}(\sigma_{t.ar>1000}(Gyarto \bowtie_{g.gyartoId=t.gyartoId} Termek)).
\]

SQL-ben ugyanez:

\begin{center}
\code{SELECT g.nev, t.nev, t.ar FROM Gyarto g INNER JOIN Termek t ON t.gyartoId = g.gyartoId WHERE t.ar > 1000;}
\end{center}

A két forma ugyanazt a logikai műveletsort írja le: összekapcsolás, szűrés, majd a szükséges mezők kiválasztása. Az SQL olvashatóbb gyakorlati parancsnyelvként, a relációs algebra pedig tömörebb formális leírásként hasznos.

\section{Gyakori vizsgahibák és összefoglalás}

\subsection{Gyakori félreértések}

\begin{itemize}
  \item \textbf{DDL és DML összekeverése.} A \code{CREATE TABLE} sémát hoz létre, az \code{INSERT} rekordot szúr be.
  \item \textbf{A \code{SELECT} nem DML-módosítás.} Lekérdez, eredményrelációt állít elő; sok felosztásban DQL-ként szerepel.
  \item \textbf{Hiányzó \code{WHERE} \code{UPDATE} vagy \code{DELETE} esetén.} Ilyenkor a művelet minden rekordra vonatkozhat.
  \item \textbf{\code{WHERE} és \code{HAVING} összekeverése.} A \code{WHERE} rekordokat szűr csoportosítás előtt, a \code{HAVING} csoportokat szűr aggregálás után.
  \item \textbf{\code{NULL} hibás vizsgálata.} \code{= NULL} helyett \code{IS NULL} kell.
  \item \textbf{Joinfeltétel elhagyása.} Ez gyakorlatilag Descartes-szorzatot vagy túl nagy eredményt okozhat.
  \item \textbf{A \code{SELECT} listában nem csoportosított mező szerepel aggregáció mellett.} Csoportos lekérdezésnél a nem aggregált mezőknek illeszkedniük kell a csoportképzéshez.
  \item \textbf{SQL és relációs algebra teljes azonosítása.} Az SQL gyakorlati nyelv, ezért \code{NULL}, ismétlődés, rendezés és dialektusbeli eltérések is vannak benne.
  \item \textbf{\code{NOT IN} használata \code{NULL}-t tartalmazó részlekérdezéssel.} Ilyenkor a \code{NOT EXISTS} sokszor biztonságosabb.
  \item \textbf{DCL elhanyagolása.} A \code{GRANT} és \code{REVOKE} nem adatmódosító, hanem jogosultságkezelő utasítás.
\end{itemize}

\subsection{Rövid vizsgaválasz-minta}

\begin{summarybox}
Az SQL relációs adatbázis-kezelő nyelv. Fő parancscsoportjai a DDL, DML, DQL és DCL. A DDL az adatbázis szerkezetét kezeli: \code{CREATE} utasítással objektumokat, például táblákat hozunk létre, \code{ALTER} utasítással módosítjuk a szerkezetüket, \code{DROP} utasítással töröljük őket. A \code{CREATE TABLE} megadja a mezőket, adattípusokat és integritási feltételeket, például elsődleges kulcsot, idegen kulcsot, \code{NOT NULL}, \code{UNIQUE}, \code{CHECK} és \code{DEFAULT} feltételeket. A DML a rekordokat kezeli: \code{INSERT} beszúr, \code{UPDATE} módosít, \code{DELETE} töröl. A DCL jogosultságkezelésre szolgál, fő utasításai a \code{GRANT} és \code{REVOKE}. A \code{SELECT} a lekérdezés központi utasítása: a \code{FROM} megadja a forrásokat, a \code{JOIN} összekapcsolja a táblákat, a \code{WHERE} rekordokat szűr, a \code{GROUP BY} csoportosít, a \code{HAVING} csoportokat szűr, a \code{SELECT} mezőlistája projekciót és számított mezőket ad, az \code{ORDER BY} pedig rendezi a megjelenített eredményt. Beágyazott \code{SELECT} szerepelhet feltételben, forrásként vagy skalár kifejezésként; lehet korrelált vagy nem korrelált. A relációs algebra műveletei SQL-ben is megvalósíthatók: a szelekció \code{WHERE}, a projekció \code{SELECT}, a join \code{JOIN}, az aggregáció \code{GROUP BY} és aggregáló függvény, a halmazműveletek pedig \code{UNION}, \code{INTERSECT} és \code{EXCEPT}/\code{MINUS} formában jelennek meg.
\end{summarybox}
